% Physical pages: 114--119 (printed pages 91--96).
% Figures: none.
% Uncertainties: none.

\chapterappendix{A}{The Symmetry Representation Theorem}

This appendix presents the proof of the fundamental theorem of
Wigner~\hyperref[ref:2.2]{[2]} that any symmetry transformation can be
represented on the Hilbert space of physical states by an operator that is
either linear and unitary or antilinear and antiunitary. For our present
purposes, the property of symmetry transformations on which we chiefly rely
is that they are ray transformations $T$ that preserve transition
probabilities, in the sense that if $\ket{\Psi_1}$ and $\ket{\Psi_2}$ are
state-vectors belonging to rays $\mathcal{R}_1$ and $\mathcal{R}_2$ then any
state-vectors $\ket{\Psi'_1}$ and $\ket{\Psi'_2}$ belonging to the
transformed rays $T\mathcal{R}_1$ and $T\mathcal{R}_2$ satisfy
\begin{equation}
  \bigl|\braket{\Psi'_1}{\Psi'_2}\bigr|^2
    =\bigl|\braket{\Psi_1}{\Psi_2}\bigr|^2.
  \tag{2.A.1}\label{eq:2.A.1}
\end{equation}
We also require that a symmetry transformation should have an inverse that
preserves transition probabilities in the same sense.

To start, consider some complete orthonormal set of state-vectors
$\ket{\Psi_k}$ belonging to rays $\mathcal{R}_k$, with
\begin{equation}
  \braket{\Psi_k}{\Psi_l}=\delta_{kl},
  \tag{2.A.2}\label{eq:2.A.2}
\end{equation}
and let $\ket{\Psi'_k}$ be some arbitrary choice of state-vectors belonging
to the transformed rays $T\mathcal{R}_k$. From Eq.~\eqref{eq:2.A.1}, we have
\begin{equation*}
  \bigl|\braket{\Psi'_k}{\Psi'_l}\bigr|^2
    =\bigl|\braket{\Psi_k}{\Psi_l}\bigr|^2
    =\delta_{kl}.
\end{equation*}
But $\braket{\Psi'_k}{\Psi'_k}$ is automatically real and positive, so
this requires that it should have the value unity, and therefore
\begin{equation}
  \braket{\Psi'_k}{\Psi'_l}=\delta_{kl}.
  \tag{2.A.3}\label{eq:2.A.3}
\end{equation}
It is easy to see that these transformed states $\ket{\Psi'_k}$ also form
a complete set, for if there were any non-zero state-vector $\ket{\Psi'}$
that was orthogonal to all of the $\ket{\Psi'_k}$, then the inverse
transform of the ray to which $\ket{\Psi'}$ belongs would consist of
non-zero state-vectors $\ket{\Psi''}$ for which, for all $k$:
\begin{equation*}
  \bigl|\braket{\Psi_k}{\Psi''}\bigr|^2
    =\bigl|\braket{\Psi'_k}{\Psi'}\bigr|^2=0,
\end{equation*}
which is impossible since the $\ket{\Psi_k}$ were assumed to form a
complete set.

We must now establish a phase convention for the states
$\ket{\Psi'_k}$. For this purpose, we single out one of the
$\ket{\Psi_k}$, say $\ket{\Psi_1}$, and consider the state-vectors
\begin{equation}
  \ket{Y_k}=\frac{1}{\sqrt{2}}
    \bigl[\ket{\Psi_1}+\ket{\Psi_k}\bigr]
  \tag{2.A.4}\label{eq:2.A.4}
\end{equation}
belonging to some ray $\mathcal{S}_k$, with $k\neq 1$. Any state-vector
$\ket{Y'_k}$ belonging to the transformed ray $T\mathcal{S}_k$ may be
expanded in the state-vectors $\ket{\Psi'_l}$,
\begin{equation*}
  \ket{Y'_k}=\sum_l c_{kl}\ket{\Psi'_l}.
\end{equation*}
From Eq.~\eqref{eq:2.A.1} we have
\begin{equation*}
  |c_{kk}|=|c_{k1}|=\frac{1}{\sqrt{2}}
\end{equation*}
and for $l\neq k$ and $l\neq 1$:
\begin{equation*}
  c_{kl}=0.
\end{equation*}
For any given $k$, by an appropriate choice of phase of the two
state-vectors $\ket{Y'_k}$ and $\ket{\Psi'_k}$, we can clearly adjust the
phases of the two non-zero coefficients $c_{kk}$ and $c_{k1}$ so that both
coefficients are just $1/\sqrt{2}$. From now on, the state-vectors
$\ket{Y'_k}$ and $\ket{\Psi'_k}$ chosen in this way will be denoted
$U\ket{Y_k}$ and $U\ket{\Psi_k}$. As we have seen,
\begin{equation}
  U\frac{1}{\sqrt{2}}
    \bigl[\ket{\Psi_k}+\ket{\Psi_1}\bigr]
  =U\ket{Y_k}
  =\frac{1}{\sqrt{2}}
    \bigl[U\ket{\Psi_k}+U\ket{\Psi_1}\bigr].
  \tag{2.A.5}\label{eq:2.A.5}
\end{equation}
However, it still remains to define $U\ket{\Psi}$ for general
state-vectors $\ket{\Psi}$.

Now consider an arbitrary state-vector $\ket{\Psi}$ belonging to an
arbitrary ray $\mathcal{R}$, and expand it in the $\ket{\Psi_k}$:
\begin{equation}
  \ket{\Psi}=\sum_k C_k\ket{\Psi_k}.
  \tag{2.A.6}\label{eq:2.A.6}
\end{equation}
Any state $\ket{\Psi'}$ that belongs to the transformed ray
$T\mathcal{R}$ may similarly be expanded in the complete orthonormal set
$U\ket{\Psi_k}$:
\begin{equation}
  \ket{\Psi'}=\sum_k C'_k U\ket{\Psi_k}.
  \tag{2.A.7}\label{eq:2.A.7}
\end{equation}
The equality of $|\braket{\Psi_k}{\Psi}|^2$ and
$|\braket{U\Psi_k}{\Psi'}|^2$ tells us that for all $k$
(including $k=1$):
\begin{equation}
  |C_k|^2=|C'_k|^2,
  \tag{2.A.8}\label{eq:2.A.8}
\end{equation}
while the equality of $|\braket{Y_k}{\Psi}|^2$ and
$|\braket{UY_k}{\Psi'}|^2$ tells us that for all $k\neq 1$:
\begin{equation}
  |C_k+C_1|^2=|C'_k+C'_1|^2.
  \tag{2.A.9}\label{eq:2.A.9}
\end{equation}
The ratio of Eqs.~\eqref{eq:2.A.9} and \eqref{eq:2.A.8} yields the formula
\begin{equation}
  \operatorname{Re}(C_k/C_1)=\operatorname{Re}(C'_k/C'_1)
  \tag{2.A.10}\label{eq:2.A.10}
\end{equation}
which with Eq.~\eqref{eq:2.A.8} also requires
\begin{equation}
  \operatorname{Im}(C_k/C_1)
    =\pm\operatorname{Im}(C'_k/C'_1)
  \tag{2.A.11}\label{eq:2.A.11}
\end{equation}
and therefore either
\begin{equation}
  C_k/C_1=C'_k/C'_1,
  \tag{2.A.12}\label{eq:2.A.12}
\end{equation}
or else
\begin{equation}
  C_k/C_1=(C'_k/C'_1)^{*}.
  \tag{2.A.13}\label{eq:2.A.13}
\end{equation}
Furthermore, we can show that the same choice must be made for each $k$.
(This step in the proof was omitted by Wigner.) To see this, suppose that
for some $k$, we have $C_k/C_1=C'_k/C'_1$, while for some $l\neq k$, we
have instead $C_l/C_1=(C'_l/C'_1)^*$. Suppose also that both ratios are
complex, so that these are really different cases. (This incidentally
requires that $k\neq 1$ and $l\neq 1$, as well as $k\neq l$.) We will show
that this is impossible.

Define a state-vector
$\ket{\Phi}=\frac{1}{\sqrt{3}}[\ket{\Psi_1}+\ket{\Psi_k}+\ket{\Psi_l}]$.
Since all the ratios of the coefficients in this state-vector are real, we
must get the same ratios in any state-vector $\ket{\Phi'}$ belonging to the
transformed ray:
\begin{equation*}
  \ket{\Phi'}=\frac{\alpha}{\sqrt{3}}
    \bigl[U\ket{\Psi_1}+U\ket{\Psi_k}+U\ket{\Psi_l}\bigr],
\end{equation*}
where $\alpha$ is a phase factor with $|\alpha|=1$. But then the equality
of the transition probabilities $|\braket{\Phi}{\Psi}|$ and
$|\braket{\Phi'}{\Psi'}|$ requires that
\begin{equation*}
  \left|1+\frac{C'_k}{C'_1}+\frac{C'_l}{C'_1}\right|^2
    =\left|1+\frac{C_k}{C_1}+\frac{C_l}{C_1}\right|^2
\end{equation*}
and hence
\begin{equation*}
  \left|1+\frac{C_k}{C_1}+\frac{C_l^*}{C_1^*}\right|^2
    =\left|1+\frac{C_k}{C_1}+\frac{C_l}{C_1}\right|^2.
\end{equation*}
This is only possible if
\begin{equation*}
  \operatorname{Re}\left(\frac{C_k C_l^*}{C_1 C_1^*}\right)
    =\operatorname{Re}\left(\frac{C_k C_l}{C_1 C_1}\right)
\end{equation*}
or, in other words, if
\begin{equation*}
  \operatorname{Im}\left(\frac{C_k}{C_1}\right)
  \operatorname{Im}\left(\frac{C_l}{C_1}\right)=0.
\end{equation*}
Hence either $C_k/C_1$ or $C_l/C_1$ must be real for any pair $k,l$, in
contradiction with our assumptions. We see then that for a given symmetry
transformation $T$ applied to a given state-vector
$\sum_k C_k\ket{\Psi_k}$, we must have either Eq.~\eqref{eq:2.A.12} for all $k$, or
else Eq.~\eqref{eq:2.A.13} for all $k$.

Wigner ruled out the second possibility, Eq.~\eqref{eq:2.A.13}, because as he
showed any symmetry transformation for which this possibility is realized
would have to involve a reversal in the time coordinate, and in the proof
he presented he was considering only symmetries like rotations that do not
affect the direction of time. Here we are treating symmetries involving
time-reversal on the same basis as all other symmetries, so we will have to
consider that, for each symmetry $T$ and state-vector
$\sum_k C_k\ket{\Psi_k}$, either Eq.~\eqref{eq:2.A.12} or Eq.~\eqref{eq:2.A.13} may apply.
Depending on which of these alternatives is realized, we will now define
$U\ket{\Psi}$ to be the particular one of the state-vectors $\ket{\Psi'}$
belonging to the ray $T\mathcal{R}$ with phase chosen so that either
$C_1=C'_1$ or $C_1=C_1'{}^*$, respectively. Then either
\begin{equation}
  U\left(\sum_k C_k\ket{\Psi_k}\right)
    =\sum_k C_k U\ket{\Psi_k}
  \tag{2.A.14}\label{eq:2.A.14}
\end{equation}
or else
\begin{equation}
  U\left(\sum_k C_k\ket{\Psi_k}\right)
    =\sum_k C_k^* U\ket{\Psi_k}.
  \tag{2.A.15}\label{eq:2.A.15}
\end{equation}

It remains to be proved that for a given symmetry transformation, we must
make the same choice between Eqs.~\eqref{eq:2.A.14} and \eqref{eq:2.A.15} for arbitrary
values of the coefficients $C_k$. Suppose that Eq.~\eqref{eq:2.A.14} applies for a
state-vector $\sum_k A_k\ket{\Psi_k}$ while Eq.~\eqref{eq:2.A.15} applies for a
state-vector $\sum_k B_k\ket{\Psi_k}$. Then the invariance of transition
probabilities requires that
\begin{equation*}
  \left|\sum_k B_k^* A_k\right|^2
    =\left|\sum_k B_k A_k\right|^2
\end{equation*}
or equivalently
\begin{equation}
  \sum_{kl}\operatorname{Im}(A_k^*A_l)
    \operatorname{Im}(B_k^*B_l)=0.
  \tag{2.A.16}\label{eq:2.A.16}
\end{equation}
We cannot rule out the possibility that Eq.~\eqref{eq:2.A.16} may be satisfied for
a pair of state-vectors $\sum_k A_k\ket{\Psi_k}$ and
$\sum_k B_k\ket{\Psi_k}$ belonging to different rays. However, for any
pair of such state-vectors, with neither $A_k$ nor $B_k$ all of the same
phase (so that Eqs.~\eqref{eq:2.A.14} and \eqref{eq:2.A.15} are not the same), we can always
find a third state-vector $\sum_k C_k\ket{\Psi_k}$ for which\footnote{If for some pair $k,l$ both
$A_k^*A_l$ and $B_k^*B_l$ are complex, then choose all $C$'s to vanish
except for $C_k$ and $C_l$, and choose these two coefficients to have
different phases. If $A_k^*A_l$ is complex but $B_k^*B_l$ is real for some
pair $k,l$, then there must be some other pair $m,n$ (possibly with either
$m$ or $n$ but not both equal to $k$ or $l$) for which $B_m^*B_n$ is
complex. If also $A_m^*A_n$ is complex, then choose all $C$'s to vanish
except for $C_m$ and $C_n$, and choose these two coefficients to have
different phase. If $A_m^*A_n$ is real, then choose all $C$'s to vanish
except for $C_k$, $C_l$, $C_m$, and $C_n$, and choose these four
coefficients all to have different phases. The case where $B_k^*B_l$ is
complex but $A_k^*A_l$ is real is handled in just the same way.}
\begin{equation}
  \sum_{kl}\operatorname{Im}(C_k^*C_l)
    \operatorname{Im}(A_k^*A_l)\neq 0
  \tag{2.A.17}\label{eq:2.A.17}
\end{equation}
and also
\begin{equation}
  \sum_{kl}\operatorname{Im}(C_k^*C_l)
    \operatorname{Im}(B_k^*B_l)\neq 0.
  \tag{2.A.18}\label{eq:2.A.18}
\end{equation}
As we have seen, it follows from Eq.~\eqref{eq:2.A.17} that the same choice between
Eqs.~\eqref{eq:2.A.14} and \eqref{eq:2.A.15} must be made for
$\sum_k A_k\ket{\Psi_k}$ and $\sum_k C_k\ket{\Psi_k}$, and it follows from
Eq.~\eqref{eq:2.A.18} that the same choice between Eqs.~\eqref{eq:2.A.14} and \eqref{eq:2.A.15} must
be made for $\sum_k B_k\ket{\Psi_k}$ and $\sum_k C_k\ket{\Psi_k}$, so the
same choice between Eqs.~\eqref{eq:2.A.14} and \eqref{eq:2.A.15} must also be made for the
two state-vectors $\sum_k A_k\ket{\Psi_k}$ and
$\sum_k B_k\ket{\Psi_k}$ with which we started. We have thus shown that
for a given symmetry transformation $T$ either all state-vectors satisfy
Eq.~\eqref{eq:2.A.14} or else they all satisfy Eq.~\eqref{eq:2.A.15}.

It is now easy to show that as we have defined it, the quantum mechanical
operator $U$ is either linear and unitary or antilinear and antiunitary.
First, suppose that Eq.~\eqref{eq:2.A.14} is satisfied for all state-vectors
$\sum_k C_k\ket{\Psi_k}$. Any two state-vectors $\ket{\Psi}$ and
$\ket{\Phi}$ may be expanded as
\begin{equation*}
  \ket{\Psi}=\sum_k A_k\ket{\Psi_k},
  \qquad
  \ket{\Phi}=\sum_k B_k\ket{\Psi_k}
\end{equation*}
and so, using Eq.~\eqref{eq:2.A.14},
\begin{align*}
  U\bigl(\alpha\ket{\Psi}+\beta\ket{\Phi}\bigr)
  &=U\sum_k(\alpha A_k+\beta B_k)\ket{\Psi_k}
    =\sum_k(\alpha A_k+\beta B_k)U\ket{\Psi_k}\\
  &=\alpha\sum_k A_k U\ket{\Psi_k}
    +\beta\sum_k B_k U\ket{\Psi_k}.
\end{align*}
Using Eq.~\eqref{eq:2.A.14} again, this gives
\begin{equation}
  U\bigl(\alpha\ket{\Psi}+\beta\ket{\Phi}\bigr)
    =\alpha U\ket{\Psi}+\beta U\ket{\Phi},
  \tag{2.A.19}\label{eq:2.A.19}
\end{equation}
so $U$ is \emph{linear}. Also, using Eqs.~\eqref{eq:2.A.2} and \eqref{eq:2.A.3}, the scalar
product of the transformed states is
\begin{equation*}
  \braket{U\Psi}{U\Phi}
    =\sum_{kl}A_k^*B_l\braket{U\Psi_k}{U\Psi_l}
    =\sum_k A_k^*B_k,
\end{equation*}
and hence
\begin{equation}
  \braket{U\Psi}{U\Phi}=\braket{\Psi}{\Phi},
  \tag{2.A.20}\label{eq:2.A.20}
\end{equation}
so $U$ is \emph{unitary}.

The case of a symmetry that satisfies Eq.~\eqref{eq:2.A.15} for all state-vectors
may be dealt with in much the same way. The reader can probably supply the
arguments without help, but since antilinear operators may be unfamiliar,
we shall give the details here anyway. Suppose that Eq.~\eqref{eq:2.A.15} is
satisfied for all state-vectors $\sum_k C_k\ket{\Psi_k}$. Any two
state-vectors $\ket{\Psi}$ and $\ket{\Phi}$ may be expanded as before, and
so:
\begin{align*}
  U\bigl(\alpha\ket{\Psi}+\beta\ket{\Phi}\bigr)
  &=U\sum_k(\alpha A_k+\beta B_k)\ket{\Psi_k}\\
  &=\sum_k(\alpha^*A_k^*+\beta^*B_k^*)U\ket{\Psi_k}\\
  &=\alpha^*\sum_k A_k^*U\ket{\Psi_k}
    +\beta^*\sum_k B_k^*U\ket{\Psi_k}.
\end{align*}
Using Eq.~\eqref{eq:2.A.15} again, this gives
\begin{equation}
  U\bigl(\alpha\ket{\Psi}+\beta\ket{\Phi}\bigr)
    =\alpha^*U\ket{\Psi}+\beta^*U\ket{\Phi},
  \tag{2.A.21}\label{eq:2.A.21}
\end{equation}
so $U$ is \emph{antilinear}. Also, using Eqs.~\eqref{eq:2.A.2} and \eqref{eq:2.A.3}, the
scalar product of the transformed states is
\begin{equation*}
  \braket{U\Psi}{U\Phi}
    =\sum_{kl}A_k B_l^*\braket{U\Psi_k}{U\Psi_l}
    =\sum_k A_k B_k^*,
\end{equation*}
and hence
\begin{equation}
  \braket{U\Psi}{U\Phi}=\braket{\Psi}{\Phi}^{*},
  \tag{2.A.22}\label{eq:2.A.22}
\end{equation}
so $U$ is \emph{antiunitary}.
